\documentclass[../mathNotesPreamble]{subfiles}

\providecommand{\relscalefact}{1.4}
\begin{document}
\relscale{\relscalefact}
  \section{1.2: Domain and Range}
  \begin{defn*}[Interval notation]
    We can specify the domain and/or range of a function using interval notation:
    \begin{itemize}
      \item Strict inequalities (exclusive) are represented using parenthesis:
        \[x>3 \Longleftrightarrow (3,\infty) \hspace*{25mm} x<9 \Longleftrightarrow (-\infty, 9)\]
        \emph{Note}: Parenthesis are \textbf{always} used when an interval's end points are infinite
      \item Nonstrict inequalities (inclusive) are represented using square brackets:
        \[7\leq x\leq 11 \Longleftrightarrow [7,11]\]
      \item Disjoint intervals can be unioned together:
        \[x\neq -2 \Longleftrightarrow (-\infty, -2) \cup (-2,\infty)\]
      \item Any combination of the above can be used:
        \[\sbrkt{-1,1}\cup (2,3]\cup (4,5)\cup [6,\infty)\]
        \begin{center}
          \begin{tikzpicture}[declare function={
            xmin=-1.75; xmax=7.75;}]
            \begin{axis}[
              axis y line=none,
              axis x line=middle,
              axis line style={black,<->},
              xmin=xmin, xmax=xmax,
              xtick={-10,-9,...,10},
              height=30mm,
              width=0.75\linewidth,
              xlabel=,
%              ticklabel style={font=\footnotesize,inner sep=0.5pt,fill=white,opacity=1.0, text opacity=1},
              every axis plot/.append style={line width=0.95pt, color=lander_blue, samples=2, line width=1.5pt}
              ]
                \addplot[-] expression[domain=-1:1]{0};
                \addplot[-] expression[domain=2:3]{0};
                \addplot[-] expression[domain=4:5]{0};
                \addplot[->] expression[domain=6:xmax]{0};
                \addplot[holdot] coordinates{(2,0)(4,0)(5,0)};
                \addplot[soldot] coordinates{(-1,0)(1,0)(3,0)(6,0)};
            \end{axis}
          \end{tikzpicture}
        \end{center}
    \end{itemize}
  \end{defn*}
  \pagebreak

  Recall that the domain is the set of all \emph{possible} input values.

  \begin{ex*}
    Find the domain of the following functions (use interval notation):
    \hspace*{\stretch{1}}
    \href{https://www.desmos.com/calculator/8iqnmynowp}{\textcolor{blue}{\underline{Graphs}}}
  \end{ex*}
  \begin{extasks}[after-item-skip=\stretch{1}](2)
    \task $f(x)=x-3$
    \task $A=\pi r^2$
    \task $y=\sqrt{x-1}$
    \task $y=\dfrac{x+2}{x^2-4}$
    \task $y=\sqrt[3]{x-4}$
    \task $f(x)=\dfrac{\sqrt[4]{x^2+4}}{(x+1)\sqrt{x-1}}$
  \end{extasks}
  \vspace*{\stretch{1}}
  \pagebreak

  \begin{ex*}
    Identify the domain and range of the following:
    \hfill \href{https://www.desmos.com/calculator/qziglbj76s}{\textcolor{blue}{\underline{Graphs}}}
    \begin{extasks}[after-item-skip=\stretch{1}](2)
      \task $x(x+1)(x+2)^2,\quad x\geq -2$
      \task $\sin(x)$
      \task $\log(x)$
      \task $2^x,\quad x\geq -1$
    \end{extasks}
    \vspace*{\stretch{1}}
  \end{ex*}
  \pagebreak

  \begin{defn*}
    A \textbf{piecewise} function is a function with different definitions for different portions of the domain.
  \end{defn*}
  \begin{ex*}
    Rewrite the following as piecewise functions:
  \end{ex*}
  \begin{extasks}[after-item-skip=\stretch{1}](2)
    \task $\abs{x}=$
    \task $\dfrac{x}{\abs{x}}=$
    \task $\abs{x-1}+\abs{4-x}=$
  \end{extasks}
  \vspace*{\stretch{1}}
  \pagebreak

  \begin{ex*}
    Graph the following piecewise function:
    \hfill \href{https://www.desmos.com/calculator/u3ebllcipf}{\textcolor{blue}{\underline{Graphs}}}

    \[f(x)=
    \begin{cases}
      x^2,& \textnormal{if } x\leq 1\\
      3,& \textnormal{if } 1<x\leq 2\\
      x,& \textnormal{if } x>2
    \end{cases}
    \]
  \end{ex*}
\pagebreak

\end{document}
